% ============================================================================
% Fichier  : partie3/P03_C09_clauses_7_10.tex
% Titre    : Support, Fonctionnement, Évaluation, Amélioration
% Auteur   : MINKA MI NGUIDJOÏ Thierry Emmanuel
% ============================================================================
% Sources : ISO/IEC 27001:2022, clauses 7 à 10. Mairie d'Eséka : cas fictif.
% ============================================================================

\chapter{Support, Fonctionnement, Évaluation, Amélioration}
\label{chap:clauses710}
\chaptermark{Clauses 7 à 10}

\epigraph{%
    « Un système se juge à ce qu'il produit comme preuves,\\
    pas à ce qu'il déclare comme intentions. »%
}{Principe du cours}

\niveauChapitre{Intermédiaire}{Le Chapitre~\ref{chap:clauses46}.}

\begin{boxobjectifs}
\begin{itemize}
    \item Identifier les ressources, compétences et communications exigées
          (clause~7).
    \item Comprendre l'exigence d'informations documentées.
    \item Décrire le fonctionnement et la maîtrise opérationnelle (clause~8).
    \item Expliquer l'évaluation des performances (clause~9).
    \item Caractériser l'amélioration et les non-conformités (clause~10).
\end{itemize}
\end{boxobjectifs}

\bigskip

% ============================================================================
\section*{Question directrice : comment faire vivre et prouver le SMSI ?}
% ============================================================================

Les clauses~7 à~10 font fonctionner le \sigle{SMSI}, le mesurent et
l'améliorent. Elles correspondent aux étapes Do, Check et Act du
\sigle{PDCA}. Ce chapitre les parcourt pour la mairie.

\medskip
\textbf{Plan du chapitre.}
\begin{enumerate}
    \item Le support (clause~7) (section~\ref{sec:c09-support}).
    \item Les informations documentées (section~\ref{sec:c09-doc}).
    \item Le fonctionnement (clause~8) (section~\ref{sec:c09-fonction}).
    \item L'évaluation des performances (clause~9) (section~\ref{sec:c09-eval}).
    \item L'amélioration (clause~10) (section~\ref{sec:c09-amelioration}).
\end{enumerate}

\separateur

% ============================================================================
\section{Le Support (Clause 7)}
\label{sec:c09-support}
% ============================================================================

\begin{boxdef}
La clause~7 exige de fournir les \termecle{ressources}, d'assurer les
\termecle{compétences}, la \termecle{sensibilisation}, la \termecle{communication}
et la maîtrise des \termecle{informations documentées}.
\end{boxdef}

Un \sigle{SMSI} sans moyens ni agents formés reste un document. La mairie doit
budgéter la sécurité, former ses agents (Chapitre~\ref{chap:orga}) et
communiquer ses règles.

% ============================================================================
\section{Les Informations Documentées}
\label{sec:c09-doc}
% ============================================================================

\begin{boxdef}
\termecle{Information documentée}~: document ou enregistrement que la norme
exige de tenir à jour et de maîtriser (politique, procédures, registres,
preuves d'audit). « Documenté » signifie tracé et retrouvable.
\end{boxdef}

La maîtrise documentaire est l'épine dorsale de l'audit~: ce qui n'est pas
documenté est réputé ne pas exister. Chaque exigence satisfaite laisse une
preuve écrite.

\begin{boxprocedure}
\textbf{Constatation.} La mairie sauvegarde l'état civil mais n'en conserve
aucune trace écrite (ni journal, ni test de restauration).

\medskip
\textbf{Analyse.} Même réelle, la sauvegarde est invérifiable~: sans
information documentée, l'auditeur ne peut la constater. La mesure existe sans
preuve~: en audit, elle ne compte pas.
\end{boxprocedure}

% ============================================================================
\section{Le Fonctionnement (Clause 8)}
\label{sec:c09-fonction}
% ============================================================================

La clause~8 exige de \termecle{planifier et maîtriser} les opérations, de
réaliser l'\termecle{appréciation des risques} à intervalles planifiés et de
mettre en \oe uvre le \termecle{traitement des risques}. C'est le « Do » du
cycle~: on exécute ce que la planification a décidé (parties IV et V).

% ============================================================================
\section{L'Évaluation des Performances (Clause 9)}
\label{sec:c09-eval}
% ============================================================================

\begin{tableaucours}{L{4.2cm}L{8.0cm}}{%
    \tableauHeader{Exigence} &
    \tableauHeader{Contenu}
}
Surveillance et mesure & Suivre les indicateurs de sécurité \\
Audit interne & Vérifier la conformité du SMSI \\
Revue de direction & La direction examine et décide \\
\end{tableaucours}
\vspace{-0.8em}
\begin{center}{\small\itshape Tableau~9.1. L'évaluation des performances (clause 9)}\end{center}

La clause~9 est le « Check »~: mesurer (Chapitre~\ref{chap:indicateurs}),
auditer et passer en revue de direction (Chapitre~\ref{chap:audit}). Sans
mesure, on ne sait pas si le \sigle{SMSI} fonctionne.

% ============================================================================
\section{L'Amélioration (Clause 10)}
\label{sec:c09-amelioration}
% ============================================================================

\begin{boxdef}
La clause~10 exige de traiter les \termecle{non-conformités} par des actions
correctives qui en éliminent la cause, et d'assurer l'\termecle{amélioration
continue} du \sigle{SMSI}.
\end{boxdef}

C'est le « Act »~: corriger les écarts à leur cause (Chapitre~\ref{chap:pdca})
et relancer le cycle. La boucle est bouclée~: de la planification (clause~6) à
l'amélioration (clause~10), le \sigle{SMSI} tourne.

\begin{boxmetier}{RSSI}
\textbf{Documenter pour durer.} La RSSI tient un dossier vivant~: politique,
appréciation des risques, plan de traitement, journaux, rapports d'audit,
revues de direction. Ce dossier \emph{est} la preuve que le \sigle{SMSI}
fonctionne.
\end{boxmetier}

\separateur

\begin{boxresume}
\textbf{Ce qu'il faut retenir de ce chapitre~:}
\begin{itemize}
    \item La \textbf{clause~7} fournit ressources, compétences, communication
          et informations documentées.
    \item Ce qui n'est pas \textbf{documenté} est réputé inexistant en audit.
    \item La \textbf{clause~8} exécute l'appréciation et le traitement des
          risques.
    \item La \textbf{clause~9} mesure, audite et passe en revue de direction.
    \item La \textbf{clause~10} corrige à la cause et améliore en continu.
\end{itemize}
\end{boxresume}

\begin{boxinfo}
\textbf{Points de contrôle de l'élève.} Je sais~: citer les exigences de
support~; définir une information documentée~; expliquer le fonctionnement~;
décrire l'évaluation des performances~; caractériser l'amélioration.
\end{boxinfo}

\bigskip

\begin{boxexo}[title={\ding{45}\quad Exercice~9.1. La clause 7 \niveaubase}]
Citez trois éléments de support exigés par la clause~7.
\end{boxexo}

\begin{boxexo}[title={\ding{45}\quad Exercice~9.2. Documenté \niveaubase}]
Qu'entend la norme par « information documentée »~? Donnez deux exemples.
\end{boxexo}

\begin{boxexo}[title={\ding{45}\quad Exercice~9.3. Les trois exigences \niveaubase}]
Citez les trois exigences de l'évaluation des performances (clause~9).
\end{boxexo}

\begin{boxexo}[title={\ding{45}\quad Exercice~9.4. Mesure sans preuve \niveaumoyen}]
Pourquoi une sauvegarde réelle mais non tracée ne compte-t-elle pas en audit~?
\end{boxexo}

\begin{boxexo}[title={\ding{45}\quad Exercice~9.5. PDCA et clauses \niveaumoyen}]
Reliez les clauses~7 à~10 aux étapes Do, Check et Act.
\end{boxexo}

\begin{boxexo}[title={\ding{45}\quad Exercice~9.6. Non-conformité \niveaumoyen}]
Distinguez corriger une non-conformité et en traiter la cause, par un exemple.
\end{boxexo}

\begin{boxexo}[title={\ding{45}\quad Exercice~9.7. Dossier de preuves \niveaumoyen}]
Listez cinq documents du dossier du \sigle{SMSI} de la mairie.
\end{boxexo}

\begin{boxexo}[title={\ding{45}\quad Exercice~9.8. Revue de direction \niveaumoyen}]
À quoi sert la revue de direction, et qui y participe~?
\end{boxexo}

\begin{boxexo}[title={\ding{45}\quad Exercice~9.9. Procédure documentaire \niveauexpert}]
Rédigez la procédure de maîtrise des informations documentées de la mairie
(création, validation, diffusion, archivage).
\end{boxexo}

\begin{boxexo}[title={\ding{45}\quad Exercice~9.10. Tableau clauses-preuves \niveauexpert}]
Pour chaque clause 7 à 10, indiquez une preuve documentée que la mairie devra
produire en audit.
\end{boxexo}

\bigskip

\begin{boxinfo}
\textbf{Pour aller plus loin.}
\begin{itemize}
    \item ISO/IEC 27001, clauses~7 à~10~\cite{iso27001}.
\end{itemize}
\end{boxinfo}

\begin{boxinfo}
\textbf{Transition.} Les exigences posées, reste le catalogue des mesures~: le
Chapitre~\ref{chap:iso27002} présente \norme{ISO/IEC 27002} et l'Annexe~A.
\end{boxinfo}
